\section{Introduction}
\label{sec:introduction}
\IEEEPARstart{W}{orkflow} applications express dependencies among processing
tasks and expose opportunities for parallel execution on heterogeneous
resources. Deploying such applications across cloud, edge, and device tiers
creates a placement problem in which execution speed, communication delay,
resource availability, and monetary cost interact. A low-cost node can be an
attractive execution target but a poor recovery destination. Similarly, adding
replicas can reduce the need for reexecution while increasing compute charges,
communication charges, and contention for limited deployment capacity.

An application deadline alone does not describe the disruption that a user is
willing to tolerate. Consider a workflow with deadline $D$, nominal makespan
$\cbase$, and an additional delay $R$ caused by recovery. Meeting the deadline
after a failure requires $\cbase+R\leq D$. A separate recovery service-level
agreement (SLA) may also require $R\leq R_{\mathrm{SLA}}$. Consequently, the
deployment has only
\begin{equation}
  \reff=\min\{R_{\mathrm{SLA}},D-\cbase\}
  \label{eq:intro_budget}
\end{equation}
seconds of effective recovery budget. A cost-minimizing deployment that uses
almost all of $D$ for nominal execution can violate the recovery requirement
even when the declared recovery window appears generous. Conversely, a fast
nominal schedule does not guarantee inexpensive or rapid recovery when its
tasks depend on vulnerable resources or communication paths.

This coupling motivates \emph{joint} placement and fault-tolerance decisions.
Active replication runs multiple task copies and allows a surviving copy to
complete the task. Passive protection deploys one copy and reexecutes affected
work after failure. Choosing a single strategy for every task overlooks
differences in task criticality, slack, communication demand, and replica
placement opportunities. The scheduling decision must therefore account for
both the nominal deployment and the recovery behavior induced by that
deployment.

Classical heterogeneous workflow heuristics, including HEFT and PEFT, provide
useful starting points for makespan-oriented scheduling
\cite{topcuoglu2002heft,arabnejad2014peft}. Reliability-aware and hybrid
fault-tolerant methods introduce replication or reexecution decisions
\cite{qin2023ramoma,taghinezhad2025recrd,ren2026hfwpf}. Their original objectives
and system assumptions, however, are not identical to a joint deadline and
recovery-delay constraint. A comparison must distinguish an algorithmic
design from the simulator, resource model, and failure model under which it is
evaluated. In particular, a common feasibility checker is necessary for
consistent comparisons, but sharing a checker does not by itself make
adaptations equivalent to their original algorithms.

We study a static, single-workflow deployment problem on a fixed computing
continuum. Tasks are placed under cumulative CPU and memory deployment
constraints, active replicas are separated by fault domain during
construction, and a deterministic recovery policy handles an explicitly
enumerated scenario set. We refer to this formulation as the
\emph{recovery-latency-constrained hybrid fault-tolerant scheduling problem}
(RLC-HFTSP). Its scenario model fixes failure locations and times; the
resulting feasibility claim applies to those scenarios rather than every
possible failure time in a physical system.

To solve this problem, we develop GCN-ALNS. Adaptive large neighborhood search
(ALNS) provides the combinatorial search framework
\cite{ropke2006alns,pisinger2007general}. Before a verified feasible deployment
is available, the search prioritizes estimated constraint violations. After
feasibility is established, it emphasizes cost. A graph convolutional network
(GCN) learns placement preferences from feasible solver demonstrations
\cite{kipf2017gcn}. Importantly, the GCN encodes a resource graph; workflow
structure enters through task features. Its scores guide search actions, while
the recovery verifier retains responsibility for certifying incumbents.

The contributions are as follows.
\begin{itemize}
  \item We formulate a deployment-dependent recovery budget that links nominal
  completion, a recovery SLA, and the largest additional delay over a finite
  fault-scenario set. The model specifies compute, communication, passive
  reserve, and recovery-delay cost components without double-counting replica
  compute cost.
  \item We design a feasibility-first ALNS solver with diversified
  initialization, adaptive operators, and stagnation restarts, and integrate
  supervised resource-graph guidance into construction, insertion, and
  destruction, including preference-based ranking of replica destinations.
  \item We establish an independent evaluation protocol that separates
  teacher generation, model selection, and held-out testing. It includes
  shared scenario verification, documented baseline adaptations, paired
  instance-level analysis, and versioned deployment artifacts.
\end{itemize}

The empirical questions concern whether learned guidance changes feasible
solution quality and search effort, which mechanisms account for those
changes, and how these effects vary with workflow scale and SLA tightness.
The manuscript does not assume that the complete solver must dominate every
alternative on every metric.
